\section{Asymptotic closure of the static and nonlinear interfaces}
\label{sec:terminal-modal-asymptotic-closure}

The interior folded-source theorem is uniform from $m=64$ onward, whereas
the five frontier signs have positive joint first-order limits without an
explicit cutoff.  This distinction is immaterial for the asymptotic lower
block: the latter requires only that all sufficiently large members satisfy
the static certificate.

\begin{theorem}[Asymptotic static directed-remainder certificate; proved here]
\label{thm:terminal-modal-static-asymptotic-closure}
For all sufficiently large $m$, every row of the two-step cone
\eqref{eq:terminal-modal-static-two-step-frontier}--
\eqref{eq:terminal-modal-static-two-step-cone} is strict.  Consequently the
local half-ratios \eqref{eq:terminal-modal-static-local-half-ratios} hold,
$d\ge h$ coordinatewise, and
\begin{equation}
 \norm{u_+}_{1,D}=\norm{(Ld)_+}_{1,D}<\frac{21}{80}q^3.
 \label{eq:terminal-modal-static-asymptotic-target}
\end{equation}
No explicit threshold for ``sufficiently large'' is asserted.
\end{theorem}

\begin{proof}
The replacement derivative theorem
\cref{thm:terminal-modal-static-derivative-replacement}, together with its
proved mass and base ledgers, gives the strict cone on every
$0\le j\le m-6$ for all $m\ge64$.  The positive joint first-order limits in
\cref{thm:terminal-modal-static-frontier-first-order} make each of the five
remaining expressions
\eqref{eq:terminal-modal-static-frontier-margin-1}--
\eqref{eq:terminal-modal-static-frontier-margin-r} strict for all sufficiently
large $m$.  There are only five expressions, so one common asymptotic cutoff
exists.  The exact equivalence in
\cref{lem:terminal-modal-static-two-step-ledger} now gives the full local
half-ratio family.  The induction in
\cref{lem:terminal-modal-static-tail-reduction} yields $d\ge h$, and the
same lemma then proves \eqref{eq:terminal-modal-static-asymptotic-target}.
\end{proof}

\begin{theorem}[Asymptotic projection and global-envelope regime; proved here]
\label{thm:terminal-modal-regime-asymptotic}
For all sufficiently large $m$, the conclusion of
\cref{lem:terminal-modal-regime} holds through
$K_m=\lfloor(8q)^{-1}\log(1/q)\rfloor$: every raw point is strictly positive
and the literal safe envelope is strictly unclipped.
\end{theorem}

\begin{proof}
Fix $m$ beyond the cutoff in
\cref{thm:terminal-modal-static-asymptotic-closure}.  The finite stopped
correlated theorem
\cref{thm:terminal-modal-finite-stopped-positivity} proves literal agreement,
projection inactivity, and zero safe subtraction for every $qk<3/50$.
Let $k_0$ be the first integer with $qk_0\ge3/50$, if such an integer occurs
before $K_m$.  Then $q(k_0-1)<3/50$, so agreement through time $k_0-1$ is
already known.  The late position-envelope lemma
\cref{lem:terminal-modal-late-position-envelope}, using
\eqref{eq:terminal-modal-static-asymptotic-target}, proves strict raw
positivity, strict envelope dominance, and literal agreement at time $k_0$.
Induction applies the same lemma at every later integer time through $K_m$.
If no such $k_0$ occurs, the early theorem already covers the full horizon.
Thus there is no gap at the early--late transition.
\end{proof}

\begin{theorem}[Unconditional asymptotic terminal logarithmic block]
\label{thm:terminal-modal-log-unconditional}
For all sufficiently large $m$, the named literal transported-center
execution has no terminal one-sided certificate during its first $K_m$
full-face steps.  Hence its terminal block has length
\begin{equation}
 \Omega\!\left(q^{-1}\log\frac1q\right).
 \label{eq:terminal-modal-log-unconditional}
\end{equation}
This is a statement only about the named exact-real recurrence.  It is not a
lower bound for other algorithms, local accelerated methods, nonlocal
response primitives, or alternative face schedules.
\end{theorem}

\begin{proof}
The proper-prefix chronology and the position and velocity entry profiles
are unconditional for the named family.  The preceding theorem discharges
the sole remaining regime premise in
\cref{thm:terminal-modal-conditional-log}.  Applying that theorem proves the
claim.
\end{proof}

\begin{corollary}[Charged terminal-certification ledger; named implementation]
\label{cor:terminal-modal-charged-work}
The prescribed terminal phase of the named implementation charges at least
\begin{equation}
 W_{\rm term}
 \ge \frac1{8q}\left\lfloor\frac1{8q}\log\frac1q\right\rfloor
 \ge \frac1{64q^2}\log\frac1q-\frac1{8q}
 \label{eq:terminal-modal-charged-work}
\end{equation}
before its terminal certificate can fire.  With
$\epsppr=\rho+\tau=2q/5$, this is
$\Omega(\epsppr^{-2}\log(1/\epsppr))$ prescribed certification work.
Relative to the log-free scale
$1/(\sqrt\alpha\epsppr)=5/(2q^2)$, the ledger has a logarithmic overhead.
\end{corollary}

\begin{proof}
The full path has volume $\vol(P_m)=2m=1/(8q)$, and the named implementation
literally scans that face at every terminal step.  Multiply this scan cost by
the $K_m$ steps from
\cref{thm:terminal-modal-log-unconditional} and use
$\lfloor x\rfloor\ge x-1$.  Substitution of
$\epsppr=2q/5$ gives the stated asymptotic form.
\end{proof}

\begin{remark}[Runtime and algorithmic scope]
The corollary concerns the prescribed certificate/termination work of this
literal implementation.  It is not a semantic-error lower bound: the iterate
may already be accurate before the named certificate fires.  It neither
refutes a soft-$O$ work target nor gives an eleven-resource or general local
algorithm lower bound.  In particular, exact CG on the $(m+1)$-dimensional
path terminates in at most $m+1=\Theta(q^{-1})$ matrix-vector products, so
the logarithmic terminal block does not transfer to arbitrary
exact-spectrum polynomial or Krylov methods.  Indeed, the full-face operator
has the distinct eigenvalues
\[
 \lambda_h=\frac{1+q^2-(1-q^2)\cos(h\pi/m)}2,
 \qquad 0\le h\le m.
\]
The full restricted optimum is interior because
$p^*(j)>q\chi^m-\rho>q(7/8-1/5)>0$.  Thus, with
$p(z)=\prod_{h=0}^m(1-z/\lambda_h)$ and
$s(z)=(1-p(z))/z$, set $c_\rho=D^{1/2}\bar c$ and
$Q=D^{1/2}\widehat QD^{-1/2}$ in the original symmetric coordinates; this
$Q$ is symmetric positive definite.  Spectral calculus gives $p(Q)=0$ and hence
$Qs(Q)c_\rho=c_\rho$.  The exact solution is
$x_\rho^*=s(Q)c_\rho$, where $s$ has degree $m$.
This is the polynomial reason exact CG terminates in at most $m+1$ steps.
\end{remark}

\begin{proposition}[A full-face exact-CG escape; comparison only]
After the named one-step singleton chronology discovers the full path, one
may discard its terminal state and solve the supplied full-face RPPR system
by zero-start exact CG.  In exact arithmetic this hybrid uses
$O(m^2)=O(q^{-2})$ charged work and returns the exact interior restricted
optimum.  It has no logarithmic terminal factor.
\end{proposition}

\begin{proof}
For the $m$ proper prefixes, the charged volumes are $1+2j$ for
$0\le j<m$, whose sum is $m^2$.  The note's explicit prefix optimum and
rank-one displacement formulas materialize the transported state: generate
the powers $\chi^i,\chi^{-i}$ sequentially and evaluate the formulas in one
pass.  Writing that state and scanning the gate and boundary on $U_j$ costs
$O(\vol(U_j))=O(1+2j)$, so the prefix response/transport bookkeeping is
absorbed in this sum.  The full-face path matvec costs $2m$, and the preceding
spectral argument gives at most $m+1$ exact-CG iterations, so the additional
matvec work is at most $2m(m+1)$.  Output and exact residual validation cost
$O(m)$.  Exact optimality gives zero full-face RPPR residual.  Since
$QD^{1/2}\one=\alpha D^{1/2}\one$, the standard RPPR bridge gives
$\norm{D^{-1/2}(x^0-x_\rho^*)}_\infty\le\rho=q/5$ (equality holds for this
family).  CG's intermediate vectors may
be signed, so this comparison is an internal full-face linear solve and lies
outside the named recurrence's projected-intermediate contract.
\end{proof}
